% ================================================================
% 02_Test_Section — 进阶测试帧
% 覆盖: 交换图 (tikz-cd) / 两栏布局 / 代码环境 / 复杂公式 / 表格
% ================================================================

% 注意: 含 tikz-cd 交换图的帧必须标记 [fragile] (beamer 与 tikz-cd 的 catcode 兼容要求)。
\begin{frame}[fragile]{交换图}
  范畴论中的泛性质 (universal property):

  \[
    \begin{tikzcd}
      A \arrow[r, "f"] \arrow[rd, "g"'] & B \arrow[d, "h"] \\
      & C
    \end{tikzcd}
  \]

  \pause

  当 $h$ 唯一存在时, 称 $h$ 为 $f$ 沿 $g$ 的余诱导态射。
\end{frame}

\begin{frame}{两栏布局}
  \begin{columns}[T]
    \begin{column}{0.46\textwidth}
      \begin{theorem}{费马大定理}{}
        对整数 $n \geq 3$, 方程
        \[
          x^n + y^n = z^n
        \]
        无非平凡整数解。
      \end{theorem}
    \end{column}
    \begin{column}{0.46\textwidth}
      \begin{itemize}
        \item 由 Andrew Wiles 于 1995 年证明
        \item 核心工具: 模形式与椭圆曲线
        \item 与谷山--志村猜想等价
      \end{itemize}
    \end{column}
  \end{columns}
\end{frame}

\begin{frame}[fragile]{代码环境}
  使用 \texttt{listings} 展示代码, 帧需标记为 \texttt{fragile}:

  \begin{code}[Python]
def fib(n):
    """递归计算斐波那契数列"""
    if n < 2:
        return n
    return fib(n - 1) + fib(n - 2)
  \end{code}
\end{frame}

\begin{frame}{复杂公式}
  \begin{align*}
    \det
    \begin{pmatrix}
      a & b \\
      c & d
    \end{pmatrix}
    &= ad - bc \\[6pt]
    \underbrace{\bigboxplus_{i=1}^{n} V_i}_{\text{直和}}
    &\cong V_1 \oplus V_2 \oplus \cdots \oplus V_n
  \end{align*}

  \pause

  符号库测试: $\Hom(A, B)$, $\Aut(G)$, $\GL_n(\K)$, $\closure(A)$, $\supp(f)$。
\end{frame}

\begin{frame}{表格}
  \begin{table}
    \centering
    \caption{各定理环境的配色对应}
    \begin{tabular}{lcl}
      \hline
      环境 & 类别 & 主色 \\
      \hline
      定义 & 基础陈述 & \textcolor{mainblue}{蓝} \\
      定理 & 推演结论 & \textcolor{mainred}{红} \\
      引理 & 推演结论 & \textcolor{mainorange}{橙} \\
      命题 & 推演结论 & \textcolor{mainpurple}{紫} \\
      例 & 补充说明 & \textcolor{maingreen}{绿} \\
      \hline
    \end{tabular}
  \end{table}
\end{frame}
